\documentclass[aspectratio=169,10pt]{beamer}

% Shared Beamer setup for Fall 2026 course slides.
\mode<presentation>{
  \usetheme{Madrid}
  \usecolortheme{default}
}

\definecolor{CalPolyGreen}{HTML}{154734}
\definecolor{CalPolyGold}{HTML}{BD8B13}
\definecolor{DeckInk}{HTML}{1F2933}
\definecolor{DeckMuted}{HTML}{5F6B7A}

\setbeamercolor{structure}{fg=CalPolyGreen}
\setbeamercolor{palette primary}{bg=CalPolyGreen,fg=white}
\setbeamercolor{palette secondary}{bg=DeckInk,fg=white}
\setbeamercolor{palette tertiary}{bg=CalPolyGold,fg=DeckInk}
\setbeamercolor{frametitle}{bg=CalPolyGreen,fg=white}
\setbeamercolor{title}{fg=CalPolyGreen}
\setbeamercolor{normal text}{fg=DeckInk,bg=white}
\setbeamercolor{block title}{bg=CalPolyGreen,fg=white}
\setbeamercolor{block body}{bg=black!3,fg=DeckInk}

\setbeamertemplate{navigation symbols}{}
\setbeamertemplate{footline}[frame number]
\setbeamertemplate{itemize item}{\color{CalPolyGold}$\blacktriangleright$}
\setbeamertemplate{itemize subitem}{\color{CalPolyGreen}$\bullet$}

\newcommand{\CourseNumber}{Course}
\newcommand{\CourseTitle}{Course Title}
\newcommand{\LectureNumber}{0}
\newcommand{\LectureTitle}{Lecture Title}
\newcommand{\LectureDate}{Date}
\newcommand{\CourseUnit}{Unit}

\newcommand{\coursenumber}[1]{\renewcommand{\CourseNumber}{#1}}
\newcommand{\coursetitle}[1]{\renewcommand{\CourseTitle}{#1}}
\newcommand{\lecturenumber}[1]{\renewcommand{\LectureNumber}{#1}}
\newcommand{\lecturetitle}[1]{\renewcommand{\LectureTitle}{#1}}
\newcommand{\lecturedate}[1]{\renewcommand{\LectureDate}{#1}}
\newcommand{\courseunit}[1]{\renewcommand{\CourseUnit}{#1}}

\author{Kirk Duran}
\institute{California Polytechnic State University, San Luis Obispo}
\date{\LectureDate}

\newcommand{\makedecktitle}{
  \begin{frame}[plain]
    \vspace{0.25in}
    {\usebeamerfont{title}\color{CalPolyGreen}\LectureTitle\par}
    \vspace{0.18in}
    {\large \CourseNumber{}: \CourseTitle\par}
    \vspace{0.08in}
    {\large Lecture \LectureNumber{} -- \LectureDate\par}
    \vfill
    {\small Kirk Duran \textbar{} Cal Poly, San Luis Obispo\par}
  \end{frame}
}

\newcommand{\placeholdernote}{
  \begin{block}{Placeholder}
    Replace these bullets with the final lecture material, examples, and in-class activities.
  \end{block}
}

\AtBeginSection[]{
  \begin{frame}{Roadmap}
    \tableofcontents[currentsection]
  \end{frame}
}


\pdfstringdefDisableCommands{\def\translate#1{#1}}

\graphicspath{{../assets/lecture-08/}}

% If an image file is not yet available, the slide displays a labeled
% placeholder using the intended filename. Place the matching file in
% ../assets/lecture-08/ to replace the placeholder automatically.
\newcommand{\lectureimage}[2][1.75in]{%
  \IfFileExists{../assets/lecture-08/#2}{%
    \includegraphics[width=\linewidth,height=#1,keepaspectratio]{#2}%
  }{%
    \fbox{%
      \parbox[c][#1][c]{0.94\linewidth}{%
        \centering
        \small\textbf{Image Placeholder}\\[0.08in]
        \scriptsize\texttt{\detokenize{#2}}
      }%
    }%
  }%
}

\newcommand{\lectureplaceholder}[2][1.75in]{%
  \fbox{%
    \parbox[c][#1][c]{0.94\linewidth}{%
      \centering
      \small\textbf{Image Placeholder}\\[0.08in]
      \scriptsize #2
    }%
  }%
}

\setbeamersize{text margin left=0.42in,text margin right=0.42in}

\coursenumber{CS 1001}
\coursetitle{Introduction to Computing}
\lecturenumber{8}
\lecturetitle{Functions I --- Abstraction}
\lecturedate{Fall 2026}
\courseunit{1. Computing and Program Execution}

\begin{document}

\makedecktitle

\begin{frame}{Today}
  \begin{itemize}
    \item Reconnect expressions and assignment with a data-first execution model.
    \item Use trace tables to make changing name-to-value bindings explicit.
    \item Recognize repeated calculations as opportunities for abstraction.
    \item Define functions as named, parameterized computations.
    \item Distinguish definitions, calls, parameters, and arguments precisely.
    \item Perform a first lightweight trace of a simple function call.
  \end{itemize}

  \begin{keyidea}
    Today we move from \emph{writing a calculation} to \emph{naming a reusable calculation}.
  \end{keyidea}
\end{frame}

\sectionframe{Part 1: From Data to Program State}

\begin{frame}{Start with the Data}
  \begin{columns}[T,onlytextwidth]
    \begin{column}{0.54\textwidth}
      \begin{itemize}
        \item Programs operate on values: integers, floating-point values, strings, Booleans, and other data.
        \item Expressions inspect or combine existing values to produce new values.
        \item Assignment associates a name with the value produced by an expression.
      \end{itemize}

      \begin{block}{Example}
        \texttt{length = 3}\\
        \texttt{width = 4}\\
        \texttt{area = length * width}
      \end{block}
    \end{column}

    \begin{column}{0.40\textwidth}
      % Intended visual: data values flowing through an expression into a new binding.
      \lectureimage[2.55in]{slide4.png}
    \end{column}
  \end{columns}
\end{frame}

\begin{frame}{Assignment Evaluates First, Then Binds}
  \begin{block}{Statement}
    \centering
    \Large\texttt{area = length * width}
  \end{block}

  \begin{enumerate}
    \item Retrieve the values currently associated with \texttt{length} and \texttt{width}.
    \item Evaluate \texttt{length * width}.
    \item Produce one resulting value.
    \item Bind the name \texttt{area} to that value.
  \end{enumerate}

  \begin{exampleblockcp}
    If \texttt{length -> 3} and \texttt{width -> 4}, then the right-hand side evaluates to \texttt{12}, after which \texttt{area -> 12}.
  \end{exampleblockcp}
\end{frame}

\begin{frame}{Trace Tables Make State Changes Explicit}
  \begin{center}
    \renewcommand{\arraystretch}{1.35}
    \begin{tabular}{c|l|c|l}
      \textbf{Step} & \textbf{Statement} & \textbf{Result} & \textbf{Bindings afterward} \\\hline
      1 & \texttt{x = 4} & \texttt{4} & \texttt{x -> 4} \\
      2 & \texttt{y = x + 2} & \texttt{6} & \texttt{x -> 4, y -> 6} \\
      3 & \texttt{x = y * 3} & \texttt{18} & \texttt{x -> 18, y -> 6}
    \end{tabular}
  \end{center}

  \begin{cpdefinition}
    A \textbf{trace table} records selected execution steps and the program state resulting from each step.
  \end{cpdefinition}

  \begin{keyidea}
    Tracing is prediction: determine the next value and state \emph{before} asking Python to confirm it.
  \end{keyidea}
\end{frame}

\begin{frame}{Practice: Trace a Short Sequence}
  \begin{block}{Program}
    \texttt{a = 5}\\
    \texttt{b = a + 3}\\
    \texttt{c = b * 2}\\
    \texttt{a = c - 4}
  \end{block}

  \begin{activity}
    Construct a four-row trace table. For each statement, record the right-hand-side value and all bindings after the statement executes.
  \end{activity}

  \begin{block}{Checkpoint}
    At the end, what are the values associated with \texttt{a}, \texttt{b}, and \texttt{c}? Which binding changed after it had already been created?
  \end{block}
\end{frame}

\sectionframe{Part 2: Repetition Creates the Need for Abstraction}

\begin{frame}{Three Cubes Repeat the Same Computational Structure}
  \begin{block}{Repeated calculations}
    \texttt{volume1 = length1 * width1 * height1}\\[0.08in]
    \texttt{volume2 = length2 * width2 * height2}\\[0.08in]
    \texttt{volume3 = length3 * width3 * height3}
  \end{block}

  \begin{itemize}
    \item The variable names differ because each cube has different data.
    \item The multiplication pattern is identical.
    \item Copying the formula preserves correctness only if every copy remains consistent.
  \end{itemize}

  \begin{activity}
    Circle the syntactic structure that is identical across all three statements. Underline the data that varies.
  \end{activity}
\end{frame}

\begin{frame}{Scaling Repetition Is a Design Problem}
  \begin{columns}[T,onlytextwidth]
    \begin{column}{0.55\textwidth}
      \begin{itemize}
        \item For 100 cubes, copying the same formula 100 times is technically possible.
        \item But repetition increases maintenance cost and the opportunity for inconsistent edits.
        \item More importantly, the program fails to express that all 100 calculations are instances of the \emph{same idea}.
      \end{itemize}

      \begin{alertblock}{Design signal}
        When only the data changes while the computational rule remains fixed, the repeated rule is a candidate for abstraction.
      \end{alertblock}
    \end{column}

    \begin{column}{0.39\textwidth}
      % Intended visual: many cubes feeding the same volume formula.
      \lectureimage[2.65in]{slide10.png}
    \end{column}
  \end{columns}
\end{frame}

\begin{frame}{Separate What Varies from What Remains Fixed}
  \begin{columns}[T,onlytextwidth]
    \begin{column}{0.46\textwidth}
      \begin{block}{Varies between cases}
        \centering
        \texttt{length}\\
        \texttt{width}\\
        \texttt{height}
      \end{block}
      \begin{itemize}
        \item These are the input data.
        \item Each cube supplies different values.
      \end{itemize}
    \end{column}

    \begin{column}{0.46\textwidth}
      \begin{block}{Remains invariant}
        \centering
        \texttt{length * width * height}
      \end{block}
      \begin{itemize}
        \item This is the computational rule.
        \item Every cube uses the same rule.
      \end{itemize}
    \end{column}
  \end{columns}
\end{frame}

\begin{frame}{Abstraction Names the Common Structure}
  \begin{cpdefinition}
    \textbf{Abstraction} identifies a common computational idea and gives it an interface that can be reused without rewriting the underlying computation.
  \end{cpdefinition}

  \begin{block}{From repeated formula to named computation}
    \centering
    \texttt{length * width * height}
    \qquad$\Longrightarrow$\qquad
    \texttt{volume(length, width, height)}
  \end{block}

  \begin{keyidea}
    Good abstraction is not merely shorter code: it makes the program's conceptual structure explicit.
  \end{keyidea}
\end{frame}

\begin{frame}{An Abstraction Has an Interface and an Implementation}
  \begin{columns}[T,onlytextwidth]
    \begin{column}{0.46\textwidth}
      \begin{block}{Interface}
        \begin{itemize}
          \item function name
          \item required inputs
          \item conceptual result
        \end{itemize}
      \end{block}
      \begin{exampleblockcp}
        \texttt{volume(length, width, height)}
      \end{exampleblockcp}
    \end{column}

    \begin{column}{0.46\textwidth}
      \begin{block}{Implementation}
        \begin{itemize}
          \item statements inside the function
          \item expressions used to compute the result
        \end{itemize}
      \end{block}
      \begin{exampleblockcp}
        \texttt{length * width * height}
      \end{exampleblockcp}
    \end{column}
  \end{columns}

  \begin{keyidea}
    The interface answers ``how do I use it?''; the implementation answers ``how is it computed?''
  \end{keyidea}
\end{frame}

\sectionframe{Part 3: Functions as Parameterized Computations}

\begin{frame}{A Function Is a Named, Parameterized Computation}
  \begin{columns}[T,onlytextwidth]
    \begin{column}{0.54\textwidth}
      \begin{cpdefinition}
        A \textbf{function} is a named computation whose parameters represent input values that may vary from one invocation to another.
      \end{cpdefinition}

      \begin{block}{Python}
        \texttt{def volume(length, width, height):}\\
        \hspace*{1em}\texttt{return length * width * height}
      \end{block}

      \begin{itemize}
        \item \texttt{volume} names the abstraction.
        \item The three parameters represent the varying dimensions.
        \item The body encodes the common rule.
      \end{itemize}
    \end{column}

    \begin{column}{0.40\textwidth}
      % Intended visual: three inputs entering a function box and one result leaving.
      \lectureimage[2.70in]{slide15.png}
    \end{column}
  \end{columns}
\end{frame}

\begin{frame}{Mathematical Functions Provide a Useful Starting Analogy}
  \begin{columns}[T,onlytextwidth]
    \begin{column}{0.46\textwidth}
      \begin{block}{Mathematics}
        \centering
        $f(x) = x^2$\\[0.12in]
        $f(3) = 9$
      \end{block}
    \end{column}

    \begin{column}{0.46\textwidth}
      \begin{block}{Python}
        \texttt{def square(x):}\\
        \hspace*{1em}\texttt{return x * x}\\[0.12in]
        \texttt{square(3)}
      \end{block}
    \end{column}
  \end{columns}

  \begin{itemize}
    \item A symbolic input stands for values supplied later.
    \item The same rule can be applied to many concrete inputs.
    \item Python functions additionally have execution semantics that we will formalize in Lecture 9.
  \end{itemize}
\end{frame}

\begin{frame}{Anatomy of a Python Function Definition}
  \begin{block}{Definition}
    \centering
    \Large\texttt{def volume(length, width, height):}
  \end{block}

  \begin{center}
    \renewcommand{\arraystretch}{1.35}
    \begin{tabular}{l|l}
      \textbf{Component} & \textbf{Role} \\\hline
      \texttt{def} & begins a function definition \\
      \texttt{volume} & function name \\
      \texttt{length, width, height} & parameter names \\
      \texttt{:} & begins the indented function body
    \end{tabular}
  \end{center}

  \begin{alertblock}{Syntax versus semantics}
    The syntax tells Python how the function is written; the abstraction tells the programmer what computation the function represents.
  \end{alertblock}
\end{frame}

\begin{frame}{The Function Body Encodes the Reusable Rule}
  \begin{block}{Definition}
    \texttt{def volume(length, width, height):}\\
    \hspace*{1em}\texttt{result = length * width * height}\\
    \hspace*{1em}\texttt{return result}
  \end{block}

  \begin{itemize}
    \item Indentation determines which statements belong to the function body.
    \item The body may contain multiple statements, not merely one expression.
    \item The body uses parameter names rather than hard-coded values so that the computation generalizes.
  \end{itemize}

  \begin{exampleblockcp}
    For today, read \texttt{return result} as indicating the function's computed result. Lecture 9 will formalize the execution semantics of \texttt{return}.
  \end{exampleblockcp}
\end{frame}

\begin{frame}{Parameters Name the Roles of the Inputs}
  \begin{block}{Definition}
    \texttt{def rectangle\_area(width, height):}\\
    \hspace*{1em}\texttt{return width * height}
  \end{block}

  \begin{itemize}
    \item A parameter is a name used by the function definition to describe one input position.
    \item Parameter names should communicate the semantic role of the supplied value.
    \item Parameters are placeholders for values that will be supplied when the function is called.
  \end{itemize}

  \begin{keyidea}
    Parameters belong to the abstraction's interface: they state what information the computation requires.
  \end{keyidea}
\end{frame}

\begin{frame}[shrink=5]{A Function Call Applies the Abstraction to Concrete Data}
  \begin{columns}[T,onlytextwidth]
    \begin{column}{0.50\textwidth}
      \begin{block}{Definition}
        \texttt{def volume(length, width, height):}\\
        \hspace*{1em}\texttt{return length * width * height}
      \end{block}
      \begin{block}{Call}
        \centering
        \Large\texttt{volume(2, 3, 4)}
      \end{block}
    \end{column}

    \begin{column}{0.44\textwidth}
      % Intended visual: one stored function definition with multiple arrows from call sites.
      \lectureimage[2.20in]{slide20.png}
    \end{column}
  \end{columns}

  \begin{keyidea}
    Definition creates the reusable rule; a call supplies concrete data to one use of that rule.
  \end{keyidea}
\end{frame}

\begin{frame}{One Definition Supports Many Calls}
  \begin{block}{Function}
    \texttt{def volume(length, width, height):}\\
    \hspace*{1em}\texttt{return length * width * height}
  \end{block}

  \begin{center}
    \renewcommand{\arraystretch}{1.35}
    \begin{tabular}{l|c}
      \textbf{Call} & \textbf{Conceptual result} \\\hline
      \texttt{volume(3, 4, 5)} & \texttt{60} \\
      \texttt{volume(6, 7, 8)} & \texttt{336} \\
      \texttt{volume(2, 3, 4)} & \texttt{24}
    \end{tabular}
  \end{center}

  \begin{keyidea}
    Reuse comes from keeping the computation fixed while supplying different input values.
  \end{keyidea}
\end{frame}

\sectionframe{Part 4: Parameters, Arguments, and Calls}

\begin{frame}{Parameters and Arguments Are Not Synonyms}
  \begin{columns}[T,onlytextwidth]
    \begin{column}{0.47\textwidth}
      \begin{block}{Parameters}
        \texttt{def volume(length, width, height):}
      \end{block}
      \begin{itemize}
        \item Names appearing in the definition.
        \item Describe the function's input positions.
      \end{itemize}
    \end{column}

    \begin{column}{0.47\textwidth}
      \begin{block}{Arguments}
        \texttt{volume(2, 3, 4)}
      \end{block}
      \begin{itemize}
        \item Expressions appearing at the call site.
        \item Supply data for a particular invocation.
      \end{itemize}
    \end{column}
  \end{columns}

  \begin{alertblock}{Terminology}
    Parameter: definition-side input name. Argument: call-side expression.
  \end{alertblock}
\end{frame}

\begin{frame}{Positional Arguments Correspond by Position}
  \begin{columns}[T,onlytextwidth]
    \begin{column}{0.52\textwidth}
      \begin{block}{Definition}
        \texttt{def difference(first, second):}\\
        \hspace*{1em}\texttt{return first - second}
      \end{block}

      \begin{block}{Call}
        \texttt{difference(10, 4)}
      \end{block}
    \end{column}

    \begin{column}{0.42\textwidth}
      \begin{block}{Correspondence}
        first argument \texttt{10} $\rightarrow$ \texttt{first}\\[0.10in]
        second argument \texttt{4} $\rightarrow$ \texttt{second}
      \end{block}
    \end{column}
  \end{columns}

  \begin{keyidea}
    For positional arguments, position determines which parameter receives which input value.
  \end{keyidea}
\end{frame}

\begin{frame}{Changing Argument Order Can Change the Meaning}
  \begin{columns}[T,onlytextwidth]
    \begin{column}{0.46\textwidth}
      \begin{block}{First call}
        \texttt{difference(10, 4)}\\[0.08in]
        \texttt{first -> 10}\\
        \texttt{second -> 4}\\[0.08in]
        result: \texttt{6}
      \end{block}
    \end{column}

    \begin{column}{0.46\textwidth}
      \begin{block}{Reversed call}
        \texttt{difference(4, 10)}\\[0.08in]
        \texttt{first -> 4}\\
        \texttt{second -> 10}\\[0.08in]
        result: \texttt{-6}
      \end{block}
    \end{column}
  \end{columns}

  \begin{alertblock}{Do not infer correspondence from names}
    Python does not inspect whether a caller's variable name resembles a parameter name when using positional arguments.
  \end{alertblock}
\end{frame}

\begin{frame}{Caller Names and Parameter Names Can Differ Completely}
  \begin{block}{Function}
    \texttt{def rectangle\_area(width, height):}\\
    \hspace*{1em}\texttt{return width * height}
  \end{block}

  \begin{block}{Caller}
    \texttt{w = 5}\\
    \texttt{h = 3}\\
    \texttt{result = rectangle\_area(w, h)}
  \end{block}

  \begin{center}
    \texttt{w -> 5 -> width}
    \qquad
    \texttt{h -> 3 -> height}
  \end{center}

  \begin{keyidea}
    The caller's names identify data in the caller; parameter names identify input roles inside the function abstraction.
  \end{keyidea}
\end{frame}

\sectionframe{Part 5: A First Function-Call Trace}

\begin{frame}{A Simple Call Extends the Data-Flow Model}
  \begin{columns}[T,onlytextwidth]
    \begin{column}{0.54\textwidth}
      \begin{block}{Program}
        \texttt{def double(n):}\\
        \hspace*{1em}\texttt{result = n * 2}\\
        \hspace*{1em}\texttt{return result}\\[0.12in]
        \texttt{x = 4}\\
        \texttt{answer = double(x)}
      \end{block}

      \begin{itemize}
        \item Before the call, \texttt{x -> 4}.
        \item The call uses the current value represented by \texttt{x}.
        \item For this invocation, that value corresponds to parameter \texttt{n}.
      \end{itemize}
    \end{column}

    \begin{column}{0.40\textwidth}
      % Intended visual: x -> 4 -> n -> computation -> result -> answer.
      \lectureimage[2.75in]{28.png}
    \end{column}
  \end{columns}
\end{frame}

\begin{frame}{The Parameter Represents the Supplied Input Value}
  \begin{columns}[T,onlytextwidth]
    \begin{column}{0.46\textwidth}
      \begin{block}{Caller}
        \texttt{x -> 4}
      \end{block}
      \begin{center}
        \Large $\downarrow$ value supplied
      \end{center}
    \end{column}

    \begin{column}{0.46\textwidth}
      \begin{block}{Function invocation}
        \texttt{n -> 4}
      \end{block}
      \begin{center}
        \texttt{result = n * 2}\\
        \texttt{result -> 8}
      \end{center}
    \end{column}
  \end{columns}

  \begin{keyidea}
    The important invariant is the value \texttt{4}; \texttt{x} and \texttt{n} are different names used in different roles.
  \end{keyidea}
\end{frame}

\begin{frame}{First Manual Function Trace}
  \begin{center}
    \renewcommand{\arraystretch}{1.30}
    \begin{tabular}{l|l}
      \textbf{Stage} & \textbf{Relevant bindings / values} \\\hline
      Before call & \texttt{x -> 4} \\
      Input correspondence & \texttt{n -> 4} \\
      Body calculation & \texttt{result -> 8} \\
      After computation & \texttt{answer -> 8, x -> 4}
    \end{tabular}
  \end{center}

  \begin{itemize}
    \item This trace intentionally suppresses detailed call mechanics.
    \item It is sufficient to track how input data enters the abstraction and how the resulting value is used by the caller.
  \end{itemize}

  \begin{alertblock}{Next lecture}
    We will replace this simplified trace with the precise execution model for argument evaluation, local bindings, \texttt{return}, and resuming the caller.
  \end{alertblock}
\end{frame}

\begin{frame}{Before and After: The Abstraction Clarifies Intent}
  \begin{columns}[T,onlytextwidth]
    \begin{column}{0.47\textwidth}
      \begin{block}{Repeated calculations}
        \texttt{v1 = l1 * w1 * h1}\\
        \texttt{v2 = l2 * w2 * h2}\\
        \texttt{v3 = l3 * w3 * h3}
      \end{block}
      \begin{itemize}
        \item The common idea is implicit.
      \end{itemize}
    \end{column}

    \begin{column}{0.47\textwidth}
      \begin{block}{Abstracted computation}
        \texttt{v1 = volume(l1, w1, h1)}\\
        \texttt{v2 = volume(l2, w2, h2)}\\
        \texttt{v3 = volume(l3, w3, h3)}
      \end{block}
      \begin{itemize}
        \item The operation \texttt{volume} is explicit.
      \end{itemize}
    \end{column}
  \end{columns}

  \begin{keyidea}
    Abstraction improves code by naming the concept, not merely by reducing character count.
  \end{keyidea}
\end{frame}

\begin{frame}{Practice: Design the Abstraction Before Writing the Function}
  \begin{block}{Repeated computation}
    A program repeatedly computes the cost of purchasing some quantity of an item using\\[0.08in]
    \centering
    \texttt{price * quantity}
  \end{block}

  \begin{activity}
    Determine: (1) a descriptive function name, (2) the parameters, (3) the computation that remains invariant, and (4) two example calls with different data.
  \end{activity}

  \begin{block}{Design target}
    Your function interface should make the calculation understandable before the reader inspects its body.
  \end{block}
\end{frame}

\begin{frame}{Practice: Read a Function from Its Interface}
  \begin{block}{Definition}
    \texttt{def miles\_to\_kilometers(miles):}\\
    \hspace*{1em}\texttt{return miles * 1.60934}
  \end{block}

  \begin{activity}
    Without tracing every operation, answer: What data does this function require? What conceptual transformation does it represent? What would \texttt{miles\_to\_kilometers(10)} mean?
  \end{activity}

  \begin{exampleblockcp}
    Then switch perspectives and identify the parameter name, argument value, and computation used by the implementation.
  \end{exampleblockcp}
\end{frame}

\sectionframe{Part 6: The Mental Model Going Forward}

\begin{frame}{The Prediction-First Workflow}
  \begin{center}
    \Large
    Predict
    $\longrightarrow$
    Trace
    $\longrightarrow$
    Execute
    $\longrightarrow$
    Compare
  \end{center}

  \vspace{0.18in}

  \begin{itemize}
    \item \textbf{Predict:} state what values and bindings you expect.
    \item \textbf{Trace:} justify the prediction one evaluation step at a time.
    \item \textbf{Execute:} use Python or the debugger as an observation tool.
    \item \textbf{Compare:} revise the mental model if observation disagrees with prediction.
  \end{itemize}

  \begin{keyidea}
    The debugger is most useful after you have a prediction to test.
  \end{keyidea}
\end{frame}

\begin{frame}{What We Are Intentionally Deferring}
  \begin{itemize}
    \item The exact runtime sequence when a function call begins and ends.
    \item The precise semantics of \texttt{return} as part of expression evaluation.
    \item Local execution environments and the lifetime of local bindings.
    \item Functions calling other functions and nested function-call expressions.
    \item Scope, multiple active calls, and the call stack.
  \end{itemize}

  \begin{block}{Why defer these?}
    First establish the abstraction: a named computation applied to input data. Then study the execution machinery that makes the abstraction work.
  \end{block}
\end{frame}

\begin{frame}{Bridge to Lecture 9}
  \begin{block}{Today}
    \centering
    \texttt{answer = double(x)}
  \end{block}

  \begin{itemize}
    \item We can already identify the abstraction, argument, parameter, and expected result.
    \item We have not yet fully explained what Python does between encountering the call and binding \texttt{answer}.
  \end{itemize}

  \begin{activity}
    What sequence of events must occur for \texttt{double(x)} to produce a value that assignment can bind to \texttt{answer}?
  \end{activity}

  \begin{keyidea}
    Lecture 9 answers that question by treating a value-returning function call as an expression with precise execution semantics.
  \end{keyidea}
\end{frame}

\begin{frame}{Takeaways}
  \begin{itemize}
    \item Expressions evaluate to values; assignment binds names to resulting values.
    \item Trace tables expose how those bindings change as statements execute.
    \item Repeated computational structure is a signal that an abstraction may be useful.
    \item A function names a reusable computation; parameters represent input roles that may vary.
    \item Defining a function and calling a function are distinct operations.
    \item Parameters appear in definitions; arguments appear at call sites and supply concrete data.
    \item A first function trace follows values from caller data into parameterized computation and back to the caller's result.
  \end{itemize}

  \begin{keyidea}
    Mental model: \textbf{data $\rightarrow$ expressions $\rightarrow$ bindings $\rightarrow$ abstraction $\rightarrow$ reusable computation}.
  \end{keyidea}
\end{frame}

\end{document}
